\def\EMtitlesize{24}
\begin{frame}{بخش ۲}
\begin{tikzpicture}[remember picture,overlay]
\node[fill=SLbar,draw=SLbarEdge,line width=0.8bp,rounded corners=4bp,minimum width=520bp,minimum height=170bp] (b) at ([yshift=-20bp]current page.center) {};
\node[text=SLthm,font=\bfseries\fontsize{30}{36}\selectfont] at ([yshift=36bp]b.center) {\EMdir{بخش ۲}};
\node[text=SLdef,font=\fontsize{24}{30}\selectfont] at ([yshift=-30bp]b.center) {\EMdir{معادلات دیفرانسیل با مشتقات جزئی}};
\end{tikzpicture}
\end{frame}

\section{معادلهٔ موجی یک‌بعدی و روش جداسازی متغیرها}
\def\EMfootchapter{فصل ۴ \ \textbullet\ معادلهٔ موجی یک‌بعدی و روش جداسازی متغیرها}
\def\EMtitlesize{24}
\begin{frame}{فصل ۴: معادلهٔ موجی یک‌بعدی و روش جداسازی متغیرها}
{\fontsize{20}{26}\selectfont\bfseries\color{SLdef}\EMalign مطالب این فصل\par}
\vspace{10bp}
\EMitem{\textbullet}{\dimexpr\linewidth-26bp\relax}{مقدمه}
\EMitem{\textbullet}{\dimexpr\linewidth-26bp\relax}{معادلهٔ دیفرانسیل موجی یک‌بعدی}
\EMitem{\textbullet}{\dimexpr\linewidth-26bp\relax}{حل معادلهٔ دیفرانسیل موجی یک‌بعدی}
\EMitem{\textbullet}{\dimexpr\linewidth-26bp\relax}{روش جداسازی متغیرها: مثال‌ها}
\EMitem{\textbullet}{\dimexpr\linewidth-26bp\relax}{معادلهٔ موجی یک‌بعدی با شرایط مرزی غیرصفر}
\EMitem{\textbullet}{\dimexpr\linewidth-26bp\relax}{معادلهٔ موجی یک‌بعدی با سرعت اولیهٔ صفر (امواج متحرک)}
\end{frame}

\def\EMtitlesize{24}
\begin{frame}{مقدمه}
\begin{EMbox}{def}{1}{}
\begin{uncoverenv}<2->
\EMp{معادله‌ای که در آن تابع و مشتقات وجود داشته باشند.}
\end{uncoverenv}
\end{EMbox}
\begin{EMbox}{def}{2}{}
\begin{uncoverenv}<3->
\EMp{معادله‌ای که بر حسب فقط یک متغیر مستقل و مشتقات تابع فقط بر حسب آن متغیر نوشته شده باشد.}
\end{uncoverenv}
\end{EMbox}
\begin{EMbox}{def}{3}{}
\begin{uncoverenv}<4->
\EMp{معادلهٔ دیفرانسیلی که تابع به چند متغیر مستقل وابسته بوده، مشتقات جزئی تابع بر حسب متغیرهای مختلف در معادله وجود داشته باشد.}
\end{uncoverenv}
\end{EMbox}
\begin{uncoverenv}<5->
\EMp{\textbf{معادلهٔ دیفرانسیل موج:}}
\vspace{3bp}
\EMeq{\frac{\partial^2u}{\partial t^2} =c^2\frac{\partial^2u}{\partial x^2},\qquad \frac{\partial^2u}{\partial t^2} =c^2\left(\frac{\partial^2u}{\partial x^2}+\frac{\partial^2u}{\partial y^2}\right),\qquad \frac{\partial^2u}{\partial t^2} =c^2\left(\frac{\partial^2u}{\partial x^2}+\frac{\partial^2u}{\partial y^2}+\frac{\partial^2u}{\partial z^2}\right)}
\end{uncoverenv}
\EMgap
\begin{uncoverenv}<6->
\EMp{\textbf{معادلهٔ دیفرانسیل انتقال حرارت:}}
\vspace{3bp}
\EMeq{\frac{\partial u}{\partial t} =c^2\frac{\partial^2u}{\partial x^2},\qquad \frac{\partial u}{\partial t} =c^2\left(\frac{\partial^2u}{\partial x^2}+\frac{\partial^2u}{\partial y^2}\right),\qquad \frac{\partial u}{\partial t} =c^2\left(\frac{\partial^2u}{\partial x^2}+\frac{\partial^2u}{\partial y^2}+\frac{\partial^2u}{\partial z^2}\right)}
\end{uncoverenv}
\end{frame}

\def\EMtitlesize{24}
\begin{frame}{مقدمه (ادامه)}
\begin{uncoverenv}<2->
\EMp{\textbf{معادلهٔ دیفرانسیل لاپلاس:}}
\vspace{3bp}
\EMeq{\frac{\partial^2u}{\partial x^2}+\frac{\partial^2u}{\partial y^2} =0,\qquad \frac{\partial^2u}{\partial x^2}+\frac{\partial^2u}{\partial y^2}+\frac{\partial^2u}{\partial z^2} =0,\qquad \nabla^2u =0}
\end{uncoverenv}
\end{frame}

\def\EMtitlesize{24}
\begin{frame}{معادلهٔ دیفرانسیل موجی یک‌بعدی}
\begin{minipage}[t]{500bp}\vspace{0pt}
\begin{uncoverenv}<2->
\EMp{مدلسازی یک تار مرتعش (به‌دست آوردن معادلهٔ حاکم بر \EMim{u(x,t)}).}
\end{uncoverenv}
\EMgap
\begin{uncoverenv}<3->
\EMp{\textbf{فرضیات فیزیکی:}}
\end{uncoverenv}
\EMgap
\begin{uncoverenv}<4->
\EMitem{\textbullet}{\dimexpr\linewidth-26bp\relax}{جرم به صورت یکنواخت در طول تار توزیع شده (تار همگن)، کاملاً انعطاف‌پذیر است.}
\end{uncoverenv}
\EMgap
\begin{uncoverenv}<5->
\EMitem{\textbullet}{\dimexpr\linewidth-26bp\relax}{تار چنان محکم کشیده شده است که نیروی جاذبه در آن قابل چشم‌پوشی است.}
\end{uncoverenv}
\EMgap
\begin{uncoverenv}<6->
\EMitem{\textbullet}{\dimexpr\linewidth-26bp\relax}{میزان انحراف تار از حالت تعادل بسیار کوچک است.}
\end{uncoverenv}
\EMgap
\end{minipage}
\EMside{184.3bp}{\uncover<7->{\EMfigure{1.414}{m04-string}{قطعهٔ \EMim{PQ} از تار مرتعش: کشش‌های \EMim{T_1} و \EMim{T_2} در دو انتها با زوایای \EMim{\alpha} و \EMim{\beta}}}}
\end{frame}

\def\EMtitlesize{24}
\begin{frame}{معادلهٔ دیفرانسیل موجی یک‌بعدی (ادامه)}
\begin{uncoverenv}<2->
\EMeq{T_1\cos\alpha =T_2\cos\beta =T =\text{const.}}
\end{uncoverenv}
\EMgap
\begin{uncoverenv}<3->
\EMeq{T_2\sin\beta-T_1\sin\alpha =\rho\,\Delta s\,\frac{\partial^2u}{\partial t^2}}
\end{uncoverenv}
\EMgap
\begin{uncoverenv}<4->
\EMp{که در آن \EMim{\rho} چگالی جرمی تار و \EMim{\Delta s} دیفرانسیل طول تار است. با تقسیم بر \EMim{T}:}
\vspace{3bp}
\EMeq{T\big(\tan\beta-\tan\alpha\big) =\rho\,\Delta s\,\frac{\partial^2u}{\partial t^2}\quad \Longrightarrow\quad T\left(\frac{\partial u}{\partial x}\bigg|_{(x+\Delta x,t)}-\frac{\partial u}{\partial x}\bigg|_{(x,t)}\right) =\rho\,\Delta s\,\frac{\partial^2u}{\partial t^2}}
\end{uncoverenv}
\EMgap
\begin{uncoverenv}<5->
\EMp{با استفاده از قضیهٔ مقدار میانی (\EMim{f(b)-f(a)=f'(l)(b-a)}، \EMim{l\in[a,b]}):}
\vspace{3bp}
\EMeq{T\left(\frac{\partial u}{\partial x}\bigg|_{(x+\Delta x,t)}-\frac{\partial u}{\partial x}\bigg|_{(x,t)}\right) =T\frac{\partial^2u}{\partial x^2}\bigg|_{(l,t)}\Delta x =\rho\,\Delta s\,\frac{\partial^2u}{\partial t^2},\qquad x<l<x+\Delta x}
\end{uncoverenv}
\EMgap
\begin{uncoverenv}<6->
\EMp{طبق فرض \EMim{\displaystyle\lim_{\Delta x\to0}\frac{\Delta x}{\Delta s}=1}، پس:}
\vspace{3bp}
\EMeq{T\frac{\partial^2u}{\partial x^2} =\rho\frac{\partial^2u}{\partial t^2},\qquad c^2\triangleq\frac T\rho\quad \Longrightarrow\quad \frac{\partial^2u}{\partial t^2} =c^2\frac{\partial^2u}{\partial x^2}}
\end{uncoverenv}
\end{frame}

\def\EMtitlesize{24}
\begin{frame}{معادلهٔ دیفرانسیل موجی یک‌بعدی (ادامه)}
\begin{uncoverenv}<2->
\EMp{برای حل یکتای این معادله نیاز به شرایط مرزی و شرایط اولیه داریم:}
\end{uncoverenv}
\EMgap
\begin{uncoverenv}<3->
\EMitem{\textbullet}{\dimexpr\linewidth-26bp\relax}{شرایط مرزی (\lr{Boundary Conditions}): \EMim{u(0,t)=u(L,t)=0}}
\end{uncoverenv}
\EMgap
\begin{uncoverenv}<4->
\EMitem{\textbullet}{\dimexpr\linewidth-26bp\relax}{شرایط اولیه (\lr{Initial Conditions}): \EMim{u(x,0)=f(x)} و \EMim{u_t(x,0)=g(x)}، که در آن \EMim{u_t\triangleq\dfrac{\partial u}{\partial t}}}
\end{uncoverenv}
\end{frame}

\def\EMtitlesize{24}
\begin{frame}{حل معادلهٔ دیفرانسیل موجی یک‌بعدی}
\begin{uncoverenv}<2->
\EMeq{\begin{cases}\dfrac{\partial^2u}{\partial t^2}=c^2\dfrac{\partial^2u}{\partial x^2}\\\noalign{\vskip 2mm} u(0,t)=u(L,t)=0\\ u(x,0)=f(x)\\ u_t(x,0)=g(x)\end{cases}}
\end{uncoverenv}
\EMgap
\begin{uncoverenv}<3->
\EMp{در اینجا برای حل این معادله از روش \textbf{جداسازی متغیرها} استفاده می‌کنیم. فرض:}
\vspace{3bp}
\EMeq{u(x,t) =X(x)\times T(t)}
\end{uncoverenv}
\EMgap
\begin{uncoverenv}<4->
\EMeq{\frac{\partial^2u}{\partial x^2} =X''(x)T(t),\qquad \frac{\partial^2u}{\partial t^2} =X(x)T''(t)\quad \Longrightarrow\quad X(x)T''(t) =c^2X''(x)T(t)\quad \Longrightarrow\quad \frac{T''}{c^2T} =\frac{X''}{X} =k}
\end{uncoverenv}
\EMgap
\begin{uncoverenv}<5->
\EMp{\EMim{k} باید حتماً یک عدد منفی باشد. بررسی حالت‌ها:}
\end{uncoverenv}
\end{frame}

\def\EMtitlesize{24}
\begin{frame}{حل معادلهٔ دیفرانسیل موجی یک‌بعدی (ادامه)}
\begin{uncoverenv}<2->
\EMp{\textbf{اگر \EMim{k=0}:}}
\vspace{3bp}
\EMeq{\frac{X''}{X} =0\  \Rightarrow\ X'' =0\  \Rightarrow\ X(x) =Ax+B,\qquad X(0) =0 \Rightarrow B =0,\quad X(L) =0 \Rightarrow A =0\  \Rightarrow\ X(x) =0}
\end{uncoverenv}
\EMgap
\begin{uncoverenv}<3->
\EMp{\textbf{اگر \EMim{k>0} (\EMim{k=\lambda^2}):}}
\vspace{3bp}
\EMeq{\frac{X''}{X} =\lambda^2\  \Rightarrow\ X''-\lambda^2X =0\  \Rightarrow\ X(x) =Ae^{\lambda x}+Be^{-\lambda x},\qquad X(0) =0 \Rightarrow A+B =0,\quad X(L) =0 \Rightarrow Ae^{\lambda L}+Be^{-\lambda L} \\[4pt] =0}
\end{uncoverenv}
\EMgap
\begin{uncoverenv}<4->
\EMeq{\Longrightarrow\ A =B =0\  \Rightarrow\ X(x) =0}
\end{uncoverenv}
\end{frame}

\def\EMtitlesize{24}
\begin{frame}{حل معادلهٔ دیفرانسیل موجی یک‌بعدی (ادامه)}
\begin{uncoverenv}<2->
\EMp{بنابراین \EMim{k} باید حتماً یک عدد منفی باشد (\EMim{k=-\lambda^2}):}
\vspace{3bp}
\EMeq{\frac{T''}{c^2T} =\frac{X''}{X} =-\lambda^2\quad \Longrightarrow\quad \begin{cases}X''+\lambda^2X=0\\ T''+c^2\lambda^2T=0\end{cases}}
\end{uncoverenv}
\EMgap
\begin{uncoverenv}<3->
\EMeq{X(x) =A\cos\lambda x+B\sin\lambda x,\qquad X(0) =0 \Rightarrow A =0}
\end{uncoverenv}
\EMgap
\begin{uncoverenv}<4->
\EMeq{X(L) =0\  \Rightarrow\ B\sin\lambda L =0\  \Rightarrow\ \lambda L =n\pi\ \ (n\in\mathbb{Z})\  \Rightarrow\ \lambda =\frac{n\pi}L =\lambda_n}
\end{uncoverenv}
\EMgap
\begin{uncoverenv}<5->
\EMeq{X_n(x) =B_n\sin\lambda_nx,\qquad \lambda_n =\frac{n\pi}L\ \ (n\in\mathbb{Z})}
\end{uncoverenv}
\EMgap
\begin{uncoverenv}<6->
\EMeq{T''+c^2\lambda_n^2T =0\  \Rightarrow\ T_n(t) =C_n\cos(c\lambda_nt)+D_n\sin(c\lambda_nt)}
\end{uncoverenv}
\EMgap
\begin{uncoverenv}<7->
\EMeq{u_n(x,t) =X_n(x)\times T_n(t) =B_n\sin\lambda_nx\times\big(C_n\cos(c\lambda_nt)+D_n\sin(c\lambda_nt)\big)}
\end{uncoverenv}
\end{frame}

\def\EMtitlesize{24}
\begin{frame}{حل معادلهٔ دیفرانسیل موجی یک‌بعدی (ادامه)}
\begin{uncoverenv}<2->
\EMp{با برهم‌نهی جواب‌ها:}
\vspace{3bp}
\EMeq{u(x,t) =\sum_{n=1}^{\infty}\left[a_n\cos\!\left(\frac{cn\pi}Lt\right)+b_n\sin\!\left(\frac{cn\pi}Lt\right)\right]\sin\frac{n\pi}Lx}
\end{uncoverenv}
\EMgap
\begin{uncoverenv}<3->
\EMp{\textbf{اعمال شرایط اولیه.} شرط \EMim{u(x,0)=f(x)}:}
\vspace{3bp}
\EMeq{u(x,0) =\sum_{n=1}^{\infty}a_n\sin\frac{n\pi}Lx =f(x)}
\end{uncoverenv}
\EMgap
\begin{uncoverenv}<4->
\EMp{تابع \EMim{f(x)} با یک تابع تناوبی و فرد برابر شده است، بنابراین برای محاسبهٔ ضرایب مجهول فوریه باید \EMim{f(x)} را یک تابع تناوبی (با دورهٔ تناوب \EMim{2L}) و فرد فرض کنیم. به این صورت خواهیم داشت:}
\vspace{3bp}
\EMeq{a_n =\frac2{2L}\int_{-L}^{L}f(x)\sin\!\left(\frac{2\pi}{2L}nx\right)dx =\frac1L\int_{-L}^{L}f(x)\sin\!\left(\frac{n\pi}Lx\right)dx =\frac2L\int_0^Lf(x)\sin\!\left(\frac{n\pi}Lx\right)dx}
\end{uncoverenv}
\end{frame}

\def\EMtitlesize{24}
\begin{frame}{حل معادلهٔ دیفرانسیل موجی یک‌بعدی (ادامه)}
\begin{uncoverenv}<2->
\EMp{شرط \EMim{u_t(x,0)=g(x)}:}
\vspace{3bp}
\EMeq{u_t(x,t) =\sum_{n=1}^{\infty}\left[-a_n\frac{cn\pi}L\sin\!\left(\frac{cn\pi}Lt\right)+b_n\frac{cn\pi}L\cos\!\left(\frac{cn\pi}Lt\right)\right]\sin\frac{n\pi}Lx}
\end{uncoverenv}
\EMgap
\begin{uncoverenv}<3->
\EMeq{u_t(x,0) =\sum_{n=1}^{\infty}b_n\frac{cn\pi}L\sin\frac{n\pi}Lx =g(x)\quad \Longrightarrow\quad b_n =\frac2{cn\pi}\int_0^Lg(x)\sin\!\left(\frac{n\pi}Lx\right)dx}
\end{uncoverenv}
\EMgap
\begin{EMbox}{imp}{3}{}
\begin{uncoverenv}<4->
\EMeq{u(x,t) =\sum_{n=1}^{\infty}\left[a_n\cos\!\left(\frac{cn\pi}Lt\right)+b_n\sin\!\left(\frac{cn\pi}Lt\right)\right]\sin\frac{n\pi}Lx}
\end{uncoverenv}
\EMgap
\begin{uncoverenv}<5->
\EMeq{a_n =\frac2L\int_0^Lf(x)\sin\!\left(\frac{n\pi}Lx\right)dx,\qquad b_n =\frac2{cn\pi}\int_0^Lg(x)\sin\!\left(\frac{n\pi}Lx\right)dx}
\end{uncoverenv}
\end{EMbox}
\end{frame}

\def\EMtitlesize{24}
\begin{frame}{حل معادلهٔ دیفرانسیل موجی یک‌بعدی (ادامه)}
\begin{EMbox}{ex}{1}{}
\EMp{جواب معادلهٔ موجی یک‌بعدی زیر را به‌دست آورید.}
\vspace{3bp}
\EMeq{\frac{\partial^2u}{\partial t^2} =\frac{\partial^2u}{\partial x^2},\qquad u(0,t) =u(\pi,t) =0,\qquad u(x,0) =k\sin2x,\qquad u_t(x,0) =0}
\EMgap
\begin{uncoverenv}<2->
\EMp{با توجه به اینکه سرعت اولیهٔ تار مرتعش صفر است، \EMim{g(x)=0\Rightarrow b_n=0}. بنابراین (\EMim{c=1}، \EMim{L=\pi}):}
\vspace{3bp}
\EMeq{u(x,t) =\sum_{n=1}^{\infty}a_n\cos\!\left(\frac{cn\pi}Lt\right)\sin\frac{n\pi}Lx =\sum_{n=1}^{\infty}a_n\cos nt\,\sin nx}
\end{uncoverenv}
\EMgap
\begin{uncoverenv}<3->
\EMeq{a_n =\frac2L\int_0^Lf(x)\sin\!\left(\frac{n\pi}Lx\right)dx =\frac2\pi\int_0^\pi k\sin2x\sin nx\,dx =\begin{cases}0&n\neq2\\ k&n=2\end{cases}}
\end{uncoverenv}
\EMgap
\begin{uncoverenv}<4->
\EMeq{u(x,t) =k\sin2x\cos2t}
\end{uncoverenv}
\end{EMbox}
\end{frame}

\def\EMtitlesize{24}
\begin{frame}{حل معادلهٔ دیفرانسیل موجی یک‌بعدی (ادامه)}
\begin{minipage}[t]{500bp}\vspace{0pt}
\end{minipage}
\EMside{0.0bp}{\uncover<2->{\EMfigure{1.356}{m04-mode2}{مود دوم ارتعاش تار: \EMim{u(x,t)=k\sin2x\cos2t} در لحظه‌های مختلف}}}
\end{frame}

\def\EMtitlesize{24}
\begin{frame}{حل معادلهٔ دیفرانسیل موجی یک‌بعدی (ادامه)}
\begin{EMbox}{ex}{1}{}
\EMp{جواب معادلهٔ موجی یک‌بعدی زیر را به‌دست آورید.}
\vspace{3bp}
\EMeq{\frac{\partial^2u}{\partial t^2} =\frac{\partial^2u}{\partial x^2},\qquad u(0,t) =u(\pi,t) =0,\qquad u(x,0) =f(x),\qquad u_t(x,0) =0}
\EMgap
\EMp{که در آن \EMim{f(x)} تابع مثلثی شکل زیر است: \EMim{f(x)=\dfrac kax} برای \EMim{0\le x\le a} و \EMim{f(x)=\dfrac k{a-\pi}(x-\pi)} برای \EMim{a\le x\le\pi}.}
\EMgap
\begin{uncoverenv}<2->
\EMp{\EMim{u_t(x,0)=g(x)=0\Rightarrow b_n=0} و}
\vspace{3bp}
\EMeq{u(x,t) =\sum_{n=1}^{\infty}a_n\cos nt\,\sin nx}
\end{uncoverenv}
\EMgap
\begin{uncoverenv}<3->
\EMeq{\begin{aligned}
a_n&=\frac2\pi\int_0^\pi f(x)\sin nx\,dx=\frac2\pi\int_0^a\frac ka\,x\sin nx\,dx+\frac2\pi\int_a^\pi\frac k{a-\pi}(x-\pi)\sin nx\,dx\\
&=\frac{2k}{a\pi}\left[x\left(\frac{-1}n\cos nx\right)-(1)\left(\frac{-1}{n^2}\sin nx\right)\right]_0^a+\frac{2k}{(a-\pi)\pi}\left[(x-\pi)\left(\frac{-1}n\cos nx\right)-(1)\left(\frac{-1}{n^2}\sin nx\right)\right]_a^\pi\\
&=\cdots=\frac{2k\sin na}{an^2(\pi-a)}
\end{aligned}}
\end{uncoverenv}
\end{EMbox}
\end{frame}

\def\EMtitlesize{24}
\begin{frame}{حل معادلهٔ دیفرانسیل موجی یک‌بعدی (ادامه)}
\begin{minipage}[t]{500bp}\vspace{0pt}
\begin{EMbox}{ex}{1}{}
\begin{uncoverenv}<2->
\EMeq{u(x,t) =\frac{2k}{a(\pi-a)}\sum_{n=1}^{\infty}\frac{\sin na}{n^2}\cos nt\,\sin nx}
\end{uncoverenv}
\end{EMbox}
\end{minipage}
\EMside{95.1bp}{\uncover<3->{\EMfigure{1.376}{m04-triangle}{توزیع اولیهٔ مثلثی \EMim{f(x)} با قلهٔ \EMim{k} در \EMim{x=a}}}}
\end{frame}

\def\EMtitlesize{24}
\begin{frame}{حل معادلهٔ دیفرانسیل موجی یک‌بعدی (ادامه)}
\begin{EMbox}{ex}{1}{}
\EMp{جواب معادلهٔ موجی یک‌بعدی زیر را به‌دست آورید.}
\vspace{3bp}
\EMeq{\frac{\partial^2u}{\partial t^2} =\frac{\partial^2u}{\partial x^2},\qquad u(0,t) =u(\pi,t) =0,\qquad u(x,0) =kx(\pi-x),\qquad u_t(x,0) =0}
\EMgap
\begin{uncoverenv}<2->
\EMp{\EMim{g(x)=0\Rightarrow b_n=0} و}
\vspace{3bp}
\EMeq{u(x,t) =\sum_{n=1}^{\infty}a_n\cos nt\,\sin nx}
\end{uncoverenv}
\EMgap
\begin{uncoverenv}<3->
\EMeq{a_n =\frac2\pi\int_0^\pi kx(\pi-x)\sin nx\,dx =\frac{2k}\pi\Big[(-x^2+\pi x)\Big(\frac{-1}n\cos nx\Big)-(-2x+\pi)\Big(\frac{-1}{n^2}\sin nx\Big)+(-2)\Big(\frac1{n^3}\cos nx\Big)\Big]_0^\pi \\[4pt] =\cdots =\frac{4k\big(1-(-1)^n\big)}{\pi n^3}}
\end{uncoverenv}
\EMgap
\begin{uncoverenv}<4->
\EMeq{u(x,t) =\frac{4k}\pi\sum_{n=1}^{\infty}\frac{1-(-1)^n}{n^3}\cos nt\,\sin nx}
\end{uncoverenv}
\end{EMbox}
\end{frame}

\def\EMtitlesize{24}
\begin{frame}{حل معادلهٔ دیفرانسیل موجی یک‌بعدی (ادامه)}
\begin{minipage}[t]{500bp}\vspace{0pt}
\end{minipage}
\EMside{0.0bp}{\uncover<2->{\EMfigure{1.376}{m04-parabola}{توزیع اولیهٔ سهمی‌شکل \EMim{f(x)=kx(\pi-x)}}}}
\end{frame}

\def\EMtitlesize{24}
\begin{frame}{حل معادلهٔ دیفرانسیل موجی یک‌بعدی (ادامه)}
\begin{EMbox}{exe}{1}{}
\EMp{(الف) مقدار نسبت \EMim{\dfrac{a_1^2}{a_1^2+a_2^2+\cdots}} را به دست آورید.}
\EMgap
\begin{uncoverenv}<2->
\EMp{(ب) نسبت \EMim{\dfrac{a_{2n-1}}{a_{2n+1}}} را بر حسب \EMim{n\in\mathbb{N}} رسم کنید.}
\end{uncoverenv}
\end{EMbox}
\end{frame}

\def\EMtitlesize{24}
\begin{frame}{روش جداسازی متغیرها: مثال‌ها}
\begin{EMbox}{ex}{1}{}
\EMp{مطلوب است حل معادلهٔ دیفرانسیل مقابل به روش جداسازی متغیرها: \EMim{xu_x=yu_y}.}
\EMgap
\begin{uncoverenv}<2->
\EMeq{x\frac{\partial u}{\partial x} =y\frac{\partial u}{\partial y},\qquad u(x,y) =X(x)\cdot Y(y)}
\end{uncoverenv}
\EMgap
\begin{uncoverenv}<3->
\EMeq{xX'Y =yXY'\  \Rightarrow\ \frac{xX'}X =\frac{yY'}Y =k}
\end{uncoverenv}
\EMgap
\begin{uncoverenv}<4->
\EMeq{\frac{X'}X =\frac kx\  \Rightarrow\ \ln X =k\ln x+c_1\  \Rightarrow\ X(x) =C_1x^k,\qquad \frac{Y'}Y =\frac ky\  \Rightarrow\ \ln Y =k\ln y+c_2\  \Rightarrow\ Y(y) =C_2y^k}
\end{uncoverenv}
\EMgap
\begin{uncoverenv}<5->
\EMeq{u(x,y) =X(x)Y(y) =C_1x^kC_2y^k =C(xy)^k}
\end{uncoverenv}
\end{EMbox}
\end{frame}

\def\EMtitlesize{24}
\begin{frame}{روش جداسازی متغیرها: مثال‌ها (ادامه)}
\begin{EMbox}{ex}{1}{}
\EMp{مطلوب است حل معادلهٔ دیفرانسیل مقابل به روش جداسازی متغیرها: \EMim{u_{xx}+u_{yy}=0}.}
\EMgap
\begin{uncoverenv}<2->
\EMeq{u(x,y) =X(x)\cdot Y(y)\  \Rightarrow\ X''Y+XY'' =0\  \Rightarrow\ \frac{X''}X =-\frac{Y''}Y =k}
\end{uncoverenv}
\EMgap
\begin{uncoverenv}<3->
\EMeq{\begin{cases}X''-kX=0\ \Rightarrow\ X(x)=Ae^{\sqrt kx}+Be^{-\sqrt kx}\\ Y''+kY=0\ \Rightarrow\ Y(y)=Ce^{\sqrt{-k}\,y}+De^{-\sqrt{-k}\,y}\end{cases}}
\end{uncoverenv}
\EMgap
\begin{uncoverenv}<4->
\EMeq{u(x,y) =X(x)Y(y) =\big(Ae^{\sqrt kx}+Be^{-\sqrt kx}\big)\big(Ce^{\sqrt{-k}\,y}+De^{-\sqrt{-k}\,y}\big)}
\end{uncoverenv}
\end{EMbox}
\end{frame}

\def\EMtitlesize{24}
\begin{frame}{روش جداسازی متغیرها: مثال‌ها (ادامه)}
\begin{EMbox}{ex}{1}{}
\EMp{مطلوب است حل معادلهٔ دیفرانسیل مقابل به روش جداسازی متغیرها: \EMim{x^2u_{xy}=-3y^2u}.}
\EMgap
\begin{uncoverenv}<2->
\EMeq{x^2\frac{\partial^2u}{\partial x\partial y}+3y^2u =0,\qquad u(x,y) =X(x)\cdot Y(y)\  \Rightarrow\ x^2X'Y'+3y^2XY =0\  \Rightarrow\ \frac{x^2X'}X =\frac{-3y^2Y}{Y'} =k}
\end{uncoverenv}
\EMgap
\begin{uncoverenv}<3->
\EMeq{\frac{X'}X =\frac k{x^2}\  \Rightarrow\ \ln X =\frac{-k}x+c_1\  \Rightarrow\ X(x) =C_1e^{-k/x}}
\end{uncoverenv}
\EMgap
\begin{uncoverenv}<4->
\EMeq{\frac{Y'}Y =\frac{-3y^2}k\  \Rightarrow\ \ln Y =-\frac{y^3}k+c_2\  \Rightarrow\ Y(y) =C_2e^{-y^3/k}}
\end{uncoverenv}
\EMgap
\begin{uncoverenv}<5->
\EMeq{u(x,y) =X(x)Y(y) =C_1e^{-k/x}\times C_2e^{-y^3/k} =Ce^{-(k/x+y^3/k)}}
\end{uncoverenv}
\end{EMbox}
\end{frame}

\def\EMtitlesize{24}
\begin{frame}{معادلهٔ موجی یک‌بعدی با شرایط مرزی غیرصفر}
\begin{EMbox}{ex}{1}{}
\EMp{جواب معادلهٔ دیفرانسیل موجی یک‌بعدی زیر را به‌دست آورید.}
\vspace{3bp}
\EMeq{\frac{\partial^2u}{\partial t^2} =c^2\frac{\partial^2u}{\partial x^2},\qquad u(0,t) =P_0,\quad u(L,t) =P_1,\quad u(x,0) =f(x),\quad u_t(x,0) =g(x)}
\EMgap
\begin{uncoverenv}<2->
\EMeq{u(x,t) =w(x,t)+v(x)\quad \Longrightarrow\quad \begin{cases}w(0,t)+v(0)=P_0\\ w(L,t)+v(L)=P_1\end{cases}}
\end{uncoverenv}
\EMgap
\begin{uncoverenv}<3->
\EMeq{\frac{\partial^2w}{\partial t^2} =c^2\left(\frac{\partial^2w}{\partial x^2}+\frac{\partial^2v}{\partial x^2}\right)}
\end{uncoverenv}
\EMgap
\begin{uncoverenv}<4->
\EMp{تابع \EMim{v(x)} را چنان انتخاب می‌کنیم که مسئلهٔ \EMim{w} شرایط مرزی همگن داشته باشد:}
\vspace{3bp}
\EMeq{\begin{cases}\dfrac{\partial^2v}{\partial x^2}=0\\ v(0)=P_0\\ v(L)=P_1\end{cases}\quad \Longrightarrow\quad v(x) =\left(\frac{P_1-P_0}L\right)x+P_0}
\end{uncoverenv}
\end{EMbox}
\end{frame}

\def\EMtitlesize{24}
\begin{frame}{معادلهٔ موجی یک‌بعدی با شرایط مرزی غیرصفر (ادامه)}
\begin{EMbox}{ex}{1}{}
\begin{uncoverenv}<2->
\EMeq{\begin{cases}\dfrac{\partial^2w}{\partial t^2}=c^2\dfrac{\partial^2w}{\partial x^2}\\ w(0,t)=w(L,t)=0\end{cases}\quad \Longrightarrow\quad w(x,t) =\sum_{n=1}^{\infty}\left[A_n\cos\!\left(\frac{cn\pi}Lt\right)+B_n\sin\!\left(\frac{cn\pi}Lt\right)\right]\sin\frac{n\pi}Lx}
\end{uncoverenv}
\EMgap
\begin{uncoverenv}<3->
\EMeq{u(x,t) =\underbrace{\sum_{n=1}^{\infty}\left[A_n\cos\!\left(\frac{cn\pi}Lt\right)+B_n\sin\!\left(\frac{cn\pi}Lt\right)\right]\sin\frac{n\pi}Lx}_{w(x,t)}+\underbrace{\left(\frac{P_1-P_0}L\right)x+P_0}_{v(x)}}
\end{uncoverenv}
\EMgap
\begin{uncoverenv}<4->
\EMp{شرط اولیهٔ \EMim{u(x,0)=f(x)}:}
\vspace{3bp}
\EMeq{\sum_{n=1}^{\infty}A_n\sin\frac{n\pi}Lx =f(x)-\left(\frac{P_1-P_0}L\right)x-P_0\triangleq f_1(x)}
\end{uncoverenv}
\EMgap
\begin{uncoverenv}<5->
\EMeq{A_n =\frac2L\int_0^Lf_1(x)\sin\frac{n\pi}Lx\,dx =\frac2L\int_0^L\left[f(x)-\left(\frac{P_1-P_0}L\right)x-P_0\right]\sin\frac{n\pi}Lx\,dx}
\end{uncoverenv}
\end{EMbox}
\end{frame}

\def\EMtitlesize{24}
\begin{frame}{معادلهٔ موجی یک‌بعدی با شرایط مرزی غیرصفر (ادامه)}
\begin{EMbox}{exe}{1}{}
\EMp{\EMim{B_n} را برای حالت (الف) \EMim{g(x)=0} و (ب) \EMim{g(x)} دلخواه به دست آورید.}
\end{EMbox}
\end{frame}

\def\EMtitlesize{24}
\begin{frame}{معادلهٔ موجی یک‌بعدی با سرعت اولیهٔ صفر (امواج متحرک)}
\begin{uncoverenv}<2->
\EMp{اگر سرعت اولیه صفر باشد (\EMim{g(x)=0}) آنگاه خواهیم داشت:}
\vspace{3bp}
\EMeq{u(x,t) =\sum_{n=1}^{\infty}a_n\cos\!\left(\frac{cn\pi}Lt\right)\sin\frac{n\pi}Lx,\qquad u(x,0) =\sum_{n=1}^{\infty}a_n\sin\frac{n\pi}Lx =f(x)\equiv\tilde f(x)}
\end{uncoverenv}
\EMgap
\begin{uncoverenv}<3->
\EMp{با استفاده از \EMim{\sin\alpha\cos\beta=\frac12\big(\sin(\alpha+\beta)+\sin(\alpha-\beta)\big)}:}
\vspace{3bp}
\EMeq{\begin{aligned}
u(x,t)&=\sum_{n=1}^{\infty}a_n\frac12\left[\sin\!\left(\frac{n\pi}Lx+\frac{cn\pi}Lt\right)+\sin\!\left(\frac{n\pi}Lx-\frac{cn\pi}Lt\right)\right]\\
&=\frac12\sum_{n=1}^{\infty}a_n\left[\sin\frac{n\pi}L(x+ct)+\sin\frac{n\pi}L(x-ct)\right]\\
&=\frac12\sum_{n=1}^{\infty}a_n\sin\frac{n\pi}L(x+ct)+\frac12\sum_{n=1}^{\infty}a_n\sin\frac{n\pi}L(x-ct)=\frac12\Big[\tilde f(x+ct)+\tilde f(x-ct)\Big]
\end{aligned}}
\end{uncoverenv}
\EMgap
\begin{EMbox}{imp}{3}{}
\begin{uncoverenv}<4->
\EMeq{u(x,t) =\frac12\Big[\tilde f(x+ct)+\tilde f(x-ct)\Big]}
\end{uncoverenv}
\end{EMbox}
\end{frame}

\def\EMtitlesize{24}
\begin{frame}{معادلهٔ موجی یک‌بعدی با سرعت اولیهٔ صفر (امواج متحرک) (ادامه)}
\begin{minipage}[t]{500bp}\vspace{0pt}
\begin{EMbox}{imp}{1}{}
\begin{uncoverenv}<2->
\EMp{\EMim{\tilde f} ادامهٔ فرد و متناوب (با دورهٔ \EMim{2L}) تابع اولیهٔ \EMim{f} است؛ جواب برهم‌نهی دو موج با نیمی از دامنهٔ اولیه است که با سرعت \EMim{c} در دو جهت مخالف حرکت می‌کنند.}
\end{uncoverenv}
\end{EMbox}
\end{minipage}
\EMside{103.5bp}{\uncover<3->{\EMfigure{1.135}{m04-traveling}{ادامهٔ فرد و متناوب \EMim{\tilde f(x)} و دو موج \EMim{\frac12\tilde f(x+ct)} و \EMim{\frac12\tilde f(x-ct)} که در دو جهت مخالف حرکت می‌کنند}}}
\end{frame}

\def\EMtitlesize{24}
\begin{frame}{معادلهٔ موجی یک‌بعدی با سرعت اولیهٔ صفر (امواج متحرک) (ادامه)}
\begin{minipage}[t]{500bp}\vspace{0pt}
\begin{EMbox}{ex}{1}{}
\EMp{جواب معادلهٔ موجی یک‌بعدی زیر را با فرض سرعت اولیهٔ صفر در زمان‌های مختلف رسم کنید.}
\vspace{3bp}
\EMeq{\frac{\partial^2u}{\partial t^2} =c^2\frac{\partial^2u}{\partial x^2},\qquad u(0,t) =u(L,t) =0,\qquad u(x,0) \\[4pt] =f(x),\qquad u_t(x,0) =0}
\EMgap
\EMeq{f(x) =\begin{cases}\dfrac{2k}Lx&0<x<\dfrac L2\\\noalign{\vskip 7mm} \dfrac{2k}L(L-x)&\dfrac L2<x<L\end{cases}}
\EMgap
\begin{uncoverenv}<2->
\EMp{با رسم \EMim{\frac12\tilde f(x+ct)} و \EMim{\frac12\tilde f(x-ct)} و جمع آنها در زمان‌های \EMim{t=0}، \EMim{t=\frac{L}{5c}}، \EMim{t=\frac{2L}{5c}}، \EMim{t=\frac{L}{2c}} و \EMim{t=\frac{3L}{5c}} شکل جواب به دست می‌آید.}
\end{uncoverenv}
\end{EMbox}
\end{minipage}
\EMside{298.3bp}{\uncover<3->{\EMfigure{1.259}{m04-triangle-time}{جواب مثال برای \EMim{t=0}، \EMim{\frac{L}{5c}}، \EMim{\frac{2L}{5c}}، \EMim{\frac{L}{2c}} و \EMim{\frac{3L}{5c}}؛ در \EMim{t=\frac{L}{2c}} شکل تار صاف است}}}
\end{frame}
